\section{Nonlinear analysis}\label{sec:nonlinear}

\subsection{Nonlinear equilibrium and Newton's method}
For a geometrically or materially nonlinear structure the discrete equilibrium is
$\bm r(\vu,\lambda)=\lambda\,\fext^{\mathrm{ref}}-\fint(\vu)=\bm 0$ on the free DOFs, with load factor $\lambda$. Each element supplies $\fint$ and its consistent tangent
$\KT=\partial\fint/\partial\vu$; the base element class falls back to $\fint=\K_e\vu_e$ and $\KT=\K_e$, so linear elements run through the same drivers (and converge in one iteration).
In the load-controlled driver, $\lambda$ increases over a prescribed array of load factors (a ramp by default, or any history including unloading). \Cref{alg:newton} gives one load step.

\begin{algorithm}[t]
\caption{One load step of incremental Newton--Raphson}\label{alg:newton}
\begin{algorithmic}[1]
\Require converged $\vu_n$, new load factor $\lambda$, tolerance $\epsilon$
\State $\vu\gets\vu_n$;\quad $\mathrm{ref}\gets\max(\|\lambda(\fext^{\mathrm{ref}})_f\|,10^{-30})$
\For{$k=1,\dots,k_{\max}$}
  \State $\bm r\gets\lambda(\fext^{\mathrm{ref}})_f-\fint(\vu)_f$
  \State \textbf{if} $\|\bm r\|<\epsilon\,\mathrm{ref}$ \textbf{or} $\|\bm r\|<\epsilon$ (plus optional step-size and energy tests) \textbf{then stop}
  \State solve $\KT(\vu)_{ff}\,\delta\vu=\bm r$
  \State (optional) Armijo backtracking on $\delta\vu$
  \State $\vu\gets\vu+\delta\vu$;\quad update iteration-level internal state
\EndFor
\State commit element state variables at $\vu$ (history for plasticity)
\end{algorithmic}
\end{algorithm}

The step-size and energy tests and the Armijo line search are off by default and are configured through the shared \texttt{NewtonOptions} object. Displacement control prescribes one DOF
and records the reaction needed to hold it, which allows following paths with a limit load.

\subsection{Path following}
Near limit points the load-controlled iteration fails. \texttt{solve\_nonlinear\_arc\_length} implements the cylindrical arc-length method of Crisfield \cite{crisfield1981}
(see also \cite{riks1979}), where the load factor is an unknown and each step satisfies
\begin{equation}
  \bm r(\vu,\lambda)=\bm 0,\qquad \|\Delta\vu\|^2=\Delta\ell^2 ,
  \label{eq:arc}
\end{equation}
with $\Delta\vu$ the step increment over the free DOFs. Each iteration writes $\delta\vu=\delta\vu_r+\delta\lambda\,\delta\vu_F$ with $\delta\vu_F=\KT^{-1}\fext^{\mathrm{ref}}$,
which turns the constraint into a quadratic in $\delta\lambda$; of the two roots, the one whose increment has the larger inner product with the increment accumulated so far in the step is taken, so the path
keeps moving forward. The arc-length $\Delta\ell$ is fixed by the user and is not adapted.

\texttt{solve\_nonlinear\_koiter\_newton} replaces the linear tangent predictor with a cubic one from an asymptotic expansion of the equilibrium path. Writing $\bm w=\vu-\vu_n$ and keeping terms up to third order,
\begin{equation}
  \bm L(\bm w)+\bm Q(\bm w,\bm w)+\bm C(\bm w,\bm w,\bm w)=\Delta\lambda\,\fext^{\mathrm{ref}},\qquad \bm L=\KT(\vu_n),
  \label{eq:koiter}
\end{equation}
with $\bm Q,\bm C$ the quadratic and cubic terms of the Taylor series of $\fint$; $\bm w$ and $\Delta\lambda$ are expanded in a scalar perturbation parameter, and the resulting linear problems
are solved with the factorised tangent. The coefficients of $\bm Q$ and $\bm C$ are extracted along the perturbation direction by central finite differences of $\fint$, with a probe length proportional to the
step's displacement scale, so elements need only provide $\fint$. The implementation follows the Koiter--Newton method of Liang et al.\ \cite{liang2014koiternewton} for a single perturbation direction; a second
driver adds a second direction built from a near-critical tangent mode when one is detected (mode interaction). Newton globalisation options (step-size test, line search) are not available in the Koiter--Newton correctors.

\subsection{Elements for large rotations and strains}
\paragraph{Trusses and beams.}
\texttt{TrussTL2D} uses the Green--Lagrange strain $E=(l^2-L_0^2)/(2L_0^2)$ with current length $l$ and St~Venant--Kirchhoff stress $S=EE$ (modulus $E$), giving
$\bm f=(SA/L_0)\,\bm d$ with $\bm d$ the current relative position; the tangent is the sum of a material and a geometric part.
\texttt{Beam2DCorotational} and \texttt{Beam3DCorotational} follow the corotational idea of separating rigid-body motion from small local deformation \cite{crisfield1990corot}.
\texttt{Beam2DReissner} is a two-node shear-deformable planar beam with exact strain measures in the Simo--Reissner sense \cite{simo1985beam}, with axial, shear and curvature strains $(\varepsilon,\gamma,\kappa)$
computed from the current configuration, and a shear correction factor passed in the material tuple.

\paragraph{Plasticity.}
For small-strain J2 plasticity with linear isotropic hardening (yield stress $\sigma_y$, hardening modulus $H$), the update is the elastic-predictor, radial-return algorithm \cite{simo1985tangent,simo1998inelasticity}.
Split the trial stress into volumetric and deviatoric parts, $\bm s^{\mathrm{tr}}=2\mu\,\mathrm{dev}(\bm\varepsilon-\bm\varepsilon^p_n)$. The yield function is
$f=\|\bm s^{\mathrm{tr}}\|-\sqrt{\tfrac23}(\sigma_y+H\alpha_n)$. If $f\le0$ the step is elastic. Otherwise
\begin{equation}
  \Delta\gamma=\frac{f}{2\mu+\tfrac23H},\qquad
  \bm s=\bm s^{\mathrm{tr}}-2\mu\,\Delta\gamma\,\bm n,\quad \bm n=\frac{\bm s^{\mathrm{tr}}}{\|\bm s^{\mathrm{tr}}\|},\qquad
  \alpha=\alpha_n+\sqrt{\tfrac23}\,\Delta\gamma,
  \label{eq:j2}
\end{equation}
a closed-form radial rescaling in deviatoric space with no local iteration. The algorithmic (consistent) tangent is obtained by differentiating the discrete update \cite{simo1985tangent}. The plane-stress variant
needs a local iteration and is a separate routine, and a kinematic-hardening version exists for the 3-D element.
State variables are evaluated afresh at every Newton iteration and committed only after a converged step.
\texttt{Tet4NeoHookean} uses the compressible Neo-Hookean energy $W=\tfrac{\mu}{2}(\bar I_1-3)+\tfrac{\kappa}{2}(J-1)^2$ with $\bar I_1=J^{-2/3}\operatorname{tr}\bm C$, $J=\det\bm F$,
$\bm C=\bm F^{\!\top}\bm F$, whose second Piola--Kirchhoff stress is
$\bm S=\mu J^{-2/3}\bigl(\bm I-\tfrac13\operatorname{tr}(\bm C)\,\bm C^{-1}\bigr)+\kappa J(J-1)\bm C^{-1}$. It does not treat the incompressible limit, which requires a mixed formulation.

\subsection{Nonlinear dynamics}
Implicit nonlinear dynamics applies Newmark's method with Newton iteration in each step. With the Newmark relations for acceleration and velocity affine in the trial displacement $\bm d$,
the residual and its exact derivative are
\begin{equation}
  \bm R(\bm d)=\M\bm a(\bm d)+\C\bm v(\bm d)+\fint(\bm d)-\fext(t+\Delta t),\qquad
  \K_{\mathrm{eff}}(\bm d)=\KT(\bm d)+a_0\M+a_1\C ,
  \label{eq:nl-newmark}
\end{equation}
so the effective matrix is re-factorised every iteration. For a linear element this reduces to the linear Newmark solver of \cref{sec:analysis} and Newton converges in one step.
The explicit nonlinear integrator uses central differences on a lumped mass, so no Newton loop is needed and each step costs one internal-force evaluation.

\subsection{Contact}
\Cref{tab:contact} lists the contact formulations. Penalty contact is a unilateral spring on the penetration $\delta=\bm n\cdot\bm u-g_0$: $\bm f=k_p\max(\delta,0)\,\bm n$ with tangent $k_p\bm n\bm n^{\!\top}$ while active
\cite{wriggers2006contact}. The Lagrange multiplier version enforces $\delta=0$ exactly through one extra unknown $\lambda$ by solving the bordered Newton system
\begin{equation}
  \begin{bmatrix}\KT&\bm n\\ \bm n^{\!\top}&0\end{bmatrix}
  \begin{bmatrix}\delta\vu\\ \delta\lambda\end{bmatrix}=
  \begin{bmatrix}\fext-\fint-\lambda\bm n\\ g_0-\bm n\cdot\vu\end{bmatrix},
  \label{eq:kkt}
\end{equation}
after first checking whether the unconstrained step penetrates (active-set strategy). The augmented-Lagrangian version replaces the multiplier solve by Uzawa updates $\bar\lambda\leftarrow\max(\bar\lambda+k_p\delta,0)$
on top of a penalty Newton solve, driving the penetration to zero for any $k_p$.

\begin{table}[t]
\centering\footnotesize
\caption{Contact formulations and their scope.}\label{tab:contact}
\begin{tabular}{>{\raggedright\arraybackslash}p{4.2cm}>{\raggedright\arraybackslash}p{4.3cm}>{\raggedright\arraybackslash}p{4.5cm}}
\toprule
Formulation & Implementation & Scope \\
\midrule
Penalty gap & \texttt{GapContactPenalty} (fixed normal), \texttt{GapContactCurvedFriction} (curved obstacle, Coulomb friction) & analytic rigid obstacle \\
Node-to-segment & \texttt{NodeToSegmentContact2D}, \texttt{...Friction} & 2-D, both bodies deformable; broad-phase helper \texttt{find\_contact\_pairs\_2d} \\
Lagrange multiplier & \texttt{solve\_contact\_\allowbreak lagrange\_\allowbreak static} & one contact node, monotonic engagement \\
Augmented Lagrangian & \texttt{solve\_contact\_\allowbreak augmented\_\allowbreak lagrange\_\allowbreak static} & same scope as Lagrange \\
\bottomrule
\end{tabular}
\end{table}

The multiplier and augmented-Lagrangian solvers handle a single contact node. General multi-point contact uses the penalty elements.
